\documentclass[11pt]{article}
\usepackage[margin=1in]{geometry}
\usepackage{amsmath,amssymb}
\usepackage[T1]{fontenc}
\usepackage{lmodern}
\usepackage{hyperref}

\title{Literature Status Note: Countable Tightness of \(P(K)\) and Maharam Type}
\author{Run \texttt{fa\_banach\_001}, lane 8}
\date{2026-06-16}

\begin{document}
\maketitle

\section*{Status}

\textbf{Literature already answered, relative-consistency negative.}
The question extracted from Krupski--Plebanek \cite{KP} has a later
literature answer under Jensen's \(\diamondsuit\): Koszmider--Silber
\cite[Theorem 1]{KS} construct a compact Hausdorff space \(K\) such
that \(P(K)\) is countably tight in the weak-star topology while \(K\)
carries a Radon probability measure of uncountable Maharam type.

\section*{Source Question}

In Section 4 of \cite{KP}, after observing that tightness
\(\leq\omega_1\) of \(P(K)\) bounds Maharam type by \(\omega_1\), the
authors ask:

\begin{quote}
If \(P(K)\) has countable tightness, must every \(\mu\in P(K)\) have
countable Maharam type?
\end{quote}

This is in arXiv:1104.1451, source lines 432--446.

\section*{Supporting Answer}

Koszmider--Silber \cite{KS} explicitly study whether small weak-star
topology on spaces of measures can coexist with nonseparable measures.
Their Theorem 1 states, assuming \(\diamondsuit\), that there is a
compact Hausdorff \(K\) for which \(P(K)\) has countable tightness but
\(K\) carries a Radon measure of uncountable type. Their abstract states
the same conclusion for a Radon probability measure of uncountable
Maharam type.

Since a nonzero finite Radon measure may be normalized without changing
the Maharam type, this supplies exactly the negation of the implication
asked in \cite{KP}: in a model of \(\diamondsuit\), \(P(K)\) is
countably tight and some probability measure in \(P(K)\) has
uncountable Maharam type.

\section*{Scope and Limitations}

This is not an original result of the run and should not be counted as a
new proof or counterexample. It is a retrieval/status result: the exact
question is already answered in later literature, modulo the additional
set-theoretic axiom \(\diamondsuit\).

The answer is not an absolute ZFC construction. Koszmider--Silber note
that the extra set-theoretic assumption is necessary in the consistency
sense, because Fremlin had shown it is consistent that such compact
spaces do not exist. They also state that weakening \(\diamondsuit\) to
CH is open.

\begin{thebibliography}{9}

\bibitem{KP}
M. Krupski and G. Plebanek,
\emph{A dichotomy for the convex spaces of probability measures},
Topology Appl. 158 (2011), 2184--2190.
arXiv:1104.1451.

\bibitem{KS}
P. Koszmider and Z. Silber,
\emph{Countably tight dual ball with a nonseparable measure},
arXiv:2312.02750.

\end{thebibliography}

\end{document}
